% Source: sources/returns/S6_04_MAPS_MANIFOLDS/sections/fourfolds_D2_vertical_corrections.tex:3–278
\paragraph{Theorem D.2.50 (Proved: all vertical corrections and new simply connected fillings).}
Let $f:S\to\mathbb P^1$ be one of the twenty-one recorded
rational elliptic surfaces with Mordell--Weil group $\mathbb ZG$,
zero section $O$, and height $\langle G,G\rangle=1/d_0$.
For integers $a,b$, put $P=aG$, $Q=bG$ and
\[
 M_D=\mathcal O_S(P-O+D),\qquad
 J_D=M_D\otimes T_Q^*M_D^{-1},
\]
where $D$ is an integral divisor supported on fibre components.
Here $T_Q:S\to S$ is the actual section-translation automorphism.
Some such $D$ makes
$J_D$ the pullback of a base line bundle if and only if
$d_0\mid ab$. For every such correction the exact identity is
\[
                  J_D\simeq f^*\mathcal O_{\mathbb P^1}(ab/d_0).
                                                        \tag{D.VC1}
\]
Thus the compatible initial line bundles are not exhausted
by requiring $P$ or $Q$ narrow at every fibre.
All corrections are determined by the integral local equations
below, with precisely the invariant component orbit sums as ambiguity.
For $a,b>0$ and any nonnegative whole-fibre weights of total $ab/d_0$,
the corrected punctured family has exactly the marked monodromies
(D.GI2). This does not identify its newly permitted original caps.

There are additional explicitly marked simply connected examples.
On OS45/default, choose $(a,b)=(2,4)$ or $(4,2)$, a simple
linearization zero at any of its four $I_1$ fibres, multiplicative
Mumford fillings at every singular fibre, no extra affine boundary
translations, and ordinary gamma clutch $e=\pm1$. Use the actual
vertical corrections specified in the proof. All sixteen choices have
\[
 \begin{gathered}
 \pi_1(X)=1,\qquad H_0(X)=H_6(X)=\mathbb Z,\\
 H_2(X)=H_3(X)=\mathbb Z/b,\qquad
 H_1(X)=H_4(X)=H_5(X)=0 .
 \end{gathered}                                             \tag{D.VC7}
\]
These are new initial and marked filling data beyond the earlier
narrowness condition. Their smooth diffeomorphism classes are not
inferred from these groups.

\begin{proof}

\emph{The integral local equation.}
Index the components of a singular fibre by $C_0,\ldots,C_{v-1}$,
with $C_0$ the component meeting $O$. Let $m$ be its positive
primitive multiplicity vector, $J=(C_iC_j)$ its intersection
matrix and $\sigma$ the permutation matrix induced by $Q$.
If $P$ meets component $C_p$, the actual multidegree of
$M_D$ on this fibre is
\[
                    c_D=e_p-e_0+Jz,
                                                        \tag{D.VC2}
\]
where $z$ is the vector of coefficients of $D$ there.
All section components have multiplicity one, so $m^tc_D=0$.
The vanishing of the component degrees of $J_D$ is exactly
\[
                 (\sigma-I)Jz=-(\sigma-I)(e_p-e_0).
                                                        \tag{D.VC3}
\]
Using the inverse permutation for pullback gives the same
invariance condition and the same solution set.
The equation is integral: it may not be replaced by its
rational solution or by the saturation of its image.

For completeness, its entire ambiguity also has an exact
description. If $z,z'$ solve (D.VC3), then
$J(\sigma-I)(z-z')=0$, since $J$ commutes with the permutation.
The integral kernel of $J$ is $\mathbb Zm$, so
$(\sigma-I)(z-z')=h m$ for some integer $h$.
Apply $I+\sigma+\cdots+\sigma^{q-1}$, where $\sigma^q=I$.
Because $\sigma m=m$, the result is $0=qh m$, hence $h=0$.
Conversely every invariant vector can be added to a solution.
Thus all corrections are one particular solution plus
\[
                         \ker(\sigma-I)\subset\mathbb Z^v.
                                                        \tag{D.VC4}
\]
In particular the ambiguity consists of integral orbit sums
of components, including the whole-fibre ambiguity.

\emph{Why vanishing component degrees is sufficient globally.}
On the generic elliptic fibre, translation acts trivially
on $\operatorname{Pic}^0$, so $J_D$ is generically trivial.
It therefore has a rational section whose divisor is vertical:
one extends a generic trivialization as a rational section
on the smooth surface. If (D.VC3) holds at every fibre,
this vertical divisor has intersection zero with every
component. On each connected fibre the negative-semidefinite
intersection matrix has kernel exactly its multiplicity
vector. The divisor on that fibre is an integral multiple
of the whole fibre, since the multiplicities have gcd one.
Consequently $J_D=f^*\mathcal O(t)$ for an integer $t$.
This proves sufficiency for the actual line bundle, not
only for its rational first Chern class. Necessity follows
by restricting a pullback to the components.

We determine $t$ without suppressing a local correction.
Choose $h>0$ with $hQ$ narrow at every singular fibre.
Then $T_{hQ}$ fixes each component divisor. Telescoping
the $h$ translated copies of $J_D$ gives
\[
 \bigotimes_{j=0}^{h-1} T_Q^{j*}J_D
       =M_D\otimes T_{hQ}^*M_D^{-1}
       =\mathcal O(P-O)\otimes
                    T_{hQ}^*\mathcal O(P-O)^{-1}.          \tag{D.VC5}
\]
The left side is $f^*\mathcal O(ht)$, since $T_Q$ preserves
the base. The last expression is in the already proved
narrow-translation case of D.2.19 and D.2.46. Its degree
is the integral height pairing
$\langle P,hQ\rangle=hab/d_0$.
Pullback on base Picard groups is injective, by restriction
to the zero section. Hence $ht=hab/d_0$, proving (D.VC1)
and the necessary divisibility $d_0\mid ab$.


\emph{The actual punctured monodromies.}
Let the chosen base linearization have whole-fibre weights $w_i$.
Choose $h$ as above and iterate the actual contracting generator
$h$ times. The canonical nonzero section of $\mathcal O(D)$
off the singular fibres telescopes in this iteration, because
$T_{hQ}D=D$ as an actual divisor. Thus the iterated punctured
family is the uncorrected narrow-translation construction with
sections $P,hQ$ and weights $h w_i$. A constant nonzero scalar
in this identification changes no monodromy.

In the actual period lattice its long generator is
$\gamma_h=h\gamma$; the elliptic generators and $\delta$ are
unchanged. If $A_i$ denotes the original monodromy and $A_{i,h}$
the iterated one, put $S=\operatorname{diag}(1,1,1,h)$. Then
\[
                         A_iS=SA_{i,h}.                    \tag{D.VC8}
\]
The narrow case (D.GI2) has $q_{i,h}=h q_i$,
$z_{i,h}=h z_i$, and central entry $r_i(hz_i)+h w_i$.
Solving (D.VC8) gives exactly the unscaled matrix (D.GI2).
This proves its integrality in every corrected case as well.
The argument concerns the punctured family and its actual long
period. It does not assume a covering of a filled original or
logarithmic cap, nor suppress any subsequently chosen affine
boundary marking.

\emph{Solvability for every recorded fibre.}
For an $I_n$ fibre the component lattice of total degree
zero modulo the image of $J$ is $\mathbb Z/n$; its quotient
map is $c\mapsto\sum_i i c_i\pmod n$.
This follows by the integral successive-difference reduction
of the cycle Laplacian, whose nonzero Smith factors are
$1,\ldots,1,n$. For $n=1$ the lattice is zero; for $n=2$
the two distinct edges are both retained.
Write $q$ for the component of $Q$, $g=\gcd(n,q)$ and
$t=n/g$. A $q$-translation-invariant multidegree has one
coefficient $u_j$ on each of its $g$ orbits. Its total degree
is zero exactly when $\sum_j u_j=0$, and its component
class is $t\sum_j j u_j\pmod n$.
Thus the subgroup represented by invariant multidegrees
is exactly $t\mathbb Z/n$. The class of $e_p-e_0$ lies in it
if and only if $t\mid p$, equivalently $n\mid pq$.
This proves integral solvability of (D.VC3) for every
multiplicative fibre.

The only other types in the twenty-one profiles are
$II,III,IV,III^*,IV^*,I_1^*,I_3^*$.
The respective component-group orders are $1,2,3,2,3,4,4$.
For these finite lists, the exact marked intersection
matrices and component permutations are supplied in
\path{code/root_vertical_corrections.py}.
Their multiplicity kernels and permutation invariance are
checked there before any equation is solved. The integral
Smith calculation for every $p,q$ in each component group
gives the same criterion
\[
              \operatorname{ord}(\Phi_f)\mid pq.           \tag{D.VC6}
\]
For $III$ and $IV$ this is also the same cycle-matrix
calculation just given. The remaining star cases involve
only $4$, $9$, $16$ and $16$ pairs respectively; their
full solutions and all integer kernel generators are
retained in the output. This is an exhaustive finite
proof for those types, not a sampling of section multiples.

For every recorded profile the component of $G$ is a unit
in each nontrivial component group, and the least common
multiple of their orders is $d_0$. These finite marked
data are checked in the same executable against the
twenty-one source records. Hence (D.VC6) for $P=aG,Q=bG$
at all fibres is equivalent to $d_0\mid ab$.
Choose a Smith solution independently at each fibre and
sum the corresponding component divisors. This constructs
the required global $D$ and proves sufficiency.
All integers $a,b$ are covered because their component
classes depend only on their residues modulo $d_0$.
The supplemental check enumerates all $5696$ such residue
pairs, using $227$ distinct local component pairs, and
finds no discrepancy with (D.VC1).

\emph{A concrete correction and the remaining topology.}
At an $I_4$ fibre with $P$ and $Q$ both meeting component
two, the original multidegree is $(-1,0,1,0)$.
Taking $D=-C_0$ changes it to $(1,-1,1,-1)$, which is
invariant under translation by two. Neither section is
narrow. The full solution set is obtained by adding
integer multiples of $C_0+C_2$ and $C_1+C_3$, as the
actual Smith kernel confirms.
Globally the profiles with such additional residue pairs
are precisely OS45, OS46, OS50 and OS55, with the listed
additive variants retained. The new initial data cannot
be discarded merely because they violate the earlier OR
condition. Their original caps carry the actual corrected
multidegree (D.VC2), which can be nonprimitive and can
be permuted nontrivially by $Q$. Computing those marked
cap maps is a separate remaining part of D.2.


\emph{The sixteen new all-Mumford examples.}
The actual OS45/default data are
\[
 \begin{gathered}
 (I_8,I_1,I_1,I_1,I_1),\qquad d_0=8,\qquad
 \overline G_{I_8}=3\pmod8,\\
 g_i=(3,0),(0,0),(2,-1),(1,-2),(0,0),\\
 (B_0,\ldots,B_4)=\begin{pmatrix}1&8\\0&1\end{pmatrix},
 \begin{pmatrix}1&1\\0&1\end{pmatrix},
 \begin{pmatrix}3&4\\-1&-1\end{pmatrix},
 \begin{pmatrix}3&1\\-4&-1\end{pmatrix},
 \begin{pmatrix}0&1\\-1&2\end{pmatrix}.
 \end{gathered}                                             \tag{D.VC9}
\]
At $I_8$ one may take the following literal coefficient vectors:
\[
\begin{array}{c|c|c}
 (a,b)&(D_{C_0},\ldots,D_{C_7})&c_D\\ \hline
 (2,4)&(-1,-1,-1,0,0,0,0,0)&(0,0,1,-1,0,0,1,-1)\\
 (4,2)&(-1,0,0,1,1,1,0,0)&(1,-1,1,-1,1,-1,1,-1).
\end{array}                                                \tag{D.VC10}
\]
They are invariant respectively under translation by four and by six;
no correction is needed at the four irreducible fibres. Both height
pairings equal one. Use the punctured marking (D.VC8), one weight one
at a selected $I_1$, weight zero elsewhere, and exactly the affine
boundary data in the statement.

Each positive-weight cap is the primitive index-one triangular cap
of D.2.43. Every other cap is the rank-one Mumford cap of its actual
matrix. Thus all local filling kernels in (D.GI3) and (D.GI5) apply
without a narrowness assumption: their proofs concern these toric
monodromies and their marked primitive angular discs. The common
piece and every ordinary boundary map are the same integral bar
models as in D.2.41. No newly corrected original cap is used.

For completeness the entire finite calculation is retained in
\path{code/fourfolds_corrected_Mumford_examples.py} and its JSON output.
The input range is exactly the two pairs, four zero positions,
and two clutch signs specified above. In all sixteen cases the
unsaturated normal closure of the toric vanishing vectors is the
whole vertical lattice $\mathbb Z\langle e_1,e_2,\delta\rangle$.
All singular peripheral loops are killed by their Mumford fillings;
the remaining ordinary relation kills $\gamma$ with coefficient
$\pm1$. This proves full simple connectivity, not only $H_1=0$.

The degree-two bar presentation has $23$ generators and $70$
relations. Its exact Smith form has $22$ unit factors and last
factor $b$. The additional relative free rank in (D.GI6) is zero:
there are five component classes and the finite-filled exterior
has first Betti number five. The only Euler contribution is the
index-one positive cap, so $\chi=2$. Integral duality and universal
coefficients now give every group in (D.VC7).
These actual matrices, vertical corrections, and boundary markings
are retained for review. Invariant orbit sums added to $D$ are not
silently identified with these marked filled spaces; their induced
boundary rigidifications must also be transported.
\end{proof}
